\subsection{Finite roots and interval-soundness}
\begin{proof}[Proof of Theorem~\ref{thm:real}]
First take $3\le p\le12$.  Define exact polynomial recurrences
\[
 X_0=Y_0=0,\quad X_{j+1}=X_j^2+c,\quad
 Y_{j+1}=2X_jY_j+1.
\]
Then $X_j=x_j(c)$ and $Y_j=x'_j(c)$ identically
over $\Q[c]$.  For each rational seed in
Table~\ref{tab:seeds}, the finite checker verifies
\begin{equation}\label{eq:finite-sign}
 P_p(\sigma_p-10^{-75})>0,\quad
 P_p(\sigma_p+10^{-75})<0,\quad
 \sup_{|c-\sigma_p|\le10^{-75}}
 (Y_{6p}(c)+Y_{3p}(c))<0.
\end{equation}
The endpoints are rational, and the derivative is evaluated
on the \emph{entire} chart by outward $2^{-512}$
interval recurrences, not by sampling.  The intermediate
value theorem and strict monotonicity prove one simple root
in each open chart.  Eighty-two further sign-separated
bisections locate it in an adjacent $2^{-256}$
rational root cell; every operation on the $X_j,Y_j$
arrays and every comparison is an integer inequality.
The resulting full-hull orbit and a two-sided
fixed-point sign bracket prove all inequalities in
\eqref{eq:order1}--\eqref{eq:order2}, $A_p<0$, and the
real monotonicity.  Thus the proof of each finite assertion
is a finite exact decision procedure with explicitly supplied
input seeds, recursion, precision and required output signs.
The complete-disk properness check is proved separately below.

For $p\ge13$ use $q=4^{-p}$, $c=-2+\eps$,
$14q\le\eps\le16q$.  The positive repelling
fixed point of $f_c$ is
\[
 \beta=\frac{1+\sqrt{9-4\eps}}2,\quad
 \lambda=2\beta=4-\delta,\quad
 \delta=\frac{4\eps}{3+\sqrt{9-4\eps}}\in(0,\eps).
\]
Setting $y_j=\beta-x_j$, exact substitution gives
\begin{equation}\label{eq:cheb}
 y_{j+1}=\lambda y_j-y_j^2,\quad
 y_2=\frac83\eps-\eps^2-
 \frac{4\eps^2}{3(3+\sqrt{9-4\eps})^2},\quad
 \frac83\eps-2\eps^2<y_2<\frac83\eps.
\end{equation}
For $j\ge2$ put $y_j=2(1-\cos\theta_j)$
on the positive small-angle branch, and put
$\Theta=2^{p-2}\theta_2$ and $\kappa=\eps/q$.
The alternating Taylor inequalities for cosine,
$\sin u\ge2u/\pi$ on $[0,\pi/2]$ and
$157/50<\pi<22/7$ imply
\begin{equation}\label{eq:angle-bds}
 \frac{\kappa}{6}(1-\eps)<\Theta^2
 <\frac{\kappa}{6}(1+3\eps),\qquad
 \frac{38}{25}<\Theta<\frac{41}{25},\qquad
 \theta_j<\frac{41}{50}\quad(2\le j<p).
\end{equation}
To see the main squared-angle comparison explicitly,
$1-3\eps/4<y_2/((8/3)\eps)<1$,
$\theta_2^2\le(\pi^2/4)y_2<7\eps$,
and $1\le\theta_2^2/y_2
\le(1-7\eps/12)^{-1}\le1+\eps$.
Multiplication by $4^{p-2}$ gives the first inequality.
In the above small-angle regime the perturbation of the
exact Chebyshev doubling map satisfies
\[
 (2-\eps)\theta_j\le\theta_{j+1}\le2\theta_j.
\]
Iteration, the $1$-Lipschitz property of cosine,
and $|\beta-2|<\eps/2$ yield
\begin{equation}\label{eq:escape}
 |\theta_p-\sqrt{\kappa/6}|\le(p+6)\eps,\qquad
 |x_p-2\cos\sqrt{\kappa/6}|\le3p\eps.
\end{equation}
In particular rational endpoint comparisons using
$38/25<\sqrt{14/6}<77/50$ and
$\sqrt{16/6}>81/50$ show
$x_p(-2+14q)>1/20$ and
$x_p(-2+16q)<-1/20$.

These opposite signs alone would not give
parameter uniqueness.  Differentiating
\eqref{eq:cheb} gives, with $Y_j=dy_j/d\eps$,
\[
 Y_{j+1}=(\lambda-2y_j)Y_j+\lambda' y_j,\quad
 r_j=\frac{\eps Y_j}{y_j},\quad
 r_{j+1}=\left(1-\frac{y_j}{\lambda-y_j}\right)r_j+
 \frac{\eps\lambda'}{\lambda-y_j}.
\]
The verified rational bounds
$r_2>99/100$, $0<y_j<7/10$ for $2\le j<p$,
$\sum y_j<21/20$,
$\sum y_j/(\lambda-y_j)<7/20$ and
$p\eps<1/300000$, together with
$\prod(1-a_j)\ge1-\sum a_j$,
prove $r_p>1/2$.  Since $\theta_p>3/2$,
$y_p>9/5$, this gives
\begin{equation}\label{eq:slope}
 \frac{dx_p}{dc}<-\frac{1}{20q}.
\end{equation}
Thus the window $|x_p|\le14q$ is
crossed exactly once by the parameter.

We next normalize the first return, including its
\emph{parameter derivative}.  Write $a_j=[z^2]f_c^j$,
with $a_1=1$ and $a_{j+1}=2x_ja_j$.
The exact cosine product identity
\[
 \prod_{j=2}^{p-1}\cos(2^{j-2}\theta_2)
       =\frac{\sin\Theta}{2^{p-2}\sin\theta_2}
\]
and bounds $1-\eps\le
2x_j/(4\cos(2^{j-2}\theta_2))\le1+p\eps$
give
\begin{equation}\label{eq:A-bds}
 A_p<0,\qquad 3/20<-A_pq<17/100.
\end{equation}
Differentiation of the $x_j,a_j$ recursions,
using $|x_j|<2$ before return, gives
$|x'_j|\le4^j/3$, $|a_j|\le4^{j-1}$,
$|a'_j|\le16^{j-1}/3$.
Consequently $\tau_p'=A'_px_p+A_px'_p
>7/(1000q^2)$ on $|x_p|\le14q$,
with $\tau_p$ traversing $[-2,1/4]$ only once.
In the subchart
$J=[-1795/1000,-1785/1000]$,
one has $0<x_p<12q$.

The full complex perturbation estimate below yields
\[
 h_\tau(w)=A_pG_p(w/A_p)=w^2+\tau+e_\tau(w),\quad
 |e_\tau(w)|<88q,\quad
 |\partial_we_\tau(w)|\le256q,\quad
 |\partial_\tau e_\tau(w)|<30000q^2
 \quad(|w|\le21/10).
\]
On $|w|\le19/10$, the spatial derivative bound is
$171q$.  For $g_\tau(w)=w^2+\tau$ let
$u_j=g_\tau^j(0)$ and $\Pi(\tau)=u_6+u_3$.
On $J$, rational evaluation gives
\[
 \Pi(-1795/1000)<-1/20,\quad
 \Pi(-1785/1000)>1/20,\quad
 \Pi'(\tau)>2,\quad
 |u_j|<9/5\quad(0\le j\le6).
\]
Set $E_0=D_0=0$ and inductively
\begin{equation}\label{eq:two-errors}
 E_{j+1}=4E_j+88q,\quad
 D_{j+1}=4D_j+25(2E_j+171q)+30000q^2.
\end{equation}
For $p=13$ (and therefore for all larger $p$)
the strict rational bounds are
$E_6+E_3<1/500$ and $D_6+D_3<1/5$.
This proves that the perturbed six/three-return
expression has opposite endpoint signs
and derivative $>2-1/5>0$ throughout $J$.
It has exactly one simple zero there.

To exclude a second root in the \emph{whole signed chart},
first the quadratic model on $[-2,0]$ is trapped in
a real invariant interval; the comparison recurrence
$E_{j+1}^{\rm wide}=(21/5)E_j^{\rm wide}+88q$
gives $E_6^{\rm wide}+E_3^{\rm wide}<1/400$.
The six rational closed nonroot segments
\[
\begin{split}
 [-2,-1.989],\ [-1.988,-1.982],\
 [-1.980,-1.894],\ [-1.892,-1.795],\\
 [-1.785,-1.756],\ [-1.753,-0.01]
\end{split}
\]
are each covered by \emph{1000} rational subintervals;
the outward interval evaluation on each entire cell
proves $|\Pi|>1/400$.  In the intervening neighborhoods
$[-1.989,-1.988]$, $[-1.982,-1.980]$ and
$[-1.894,-1.892]$, the required signed normalized
itinerary is contradicted respectively by the
fourth or fifth return.  On
$[-1.756,-1.753]$ put $v=h_\tau^3(0)$:
the bounds $|v|<1/40$ and
$|(h_\tau^3)'(w)|<1$ for all $w\in[0,v]$
give
\[
 h_\tau^6(0)+h_\tau^3(0)
 =v\left(2+\int_0^1(h_\tau^3)'(sv)\,ds\right),
\]
so a zero would force $v=0$, which is forbidden by the
strict $x_{3p}<0$ condition.
For $-0.01<\tau<0$, evenness gives
$e_\tau(w)=O(w^4)$; the order-four Schwarz
lemma on $|w|\le12/5$ yields
$|e_\tau(\tau)|<|\tau|/100$ and then
$h_\tau^2(0)=\tau(1+\tau)+e_\tau(\tau)<0$,
also forbidden.  Finally $\tau\ge0$
contradicts $x_p>0$.
These analytic exclusions and all finite rational
cells, \emph{not} an unisolated broad polynomial root,
prove uniqueness in $\mathcal K_p$.

At that root the normalized model's strict
six-return cut bounds with errors \eqref{eq:two-errors}
give \eqref{eq:order1} and
\eqref{eq:order2}.  The fixed point is bracketed
on the relevant interval with strictly negative
$h_\tau'$ and opposite endpoint signs.  The
negative factor $A_p$ reverses the normalized
real cut order, which is included in the
outward computations.  Finally the derivative
of $2\cos\sqrt{\kappa/6}$ has absolute value
$>1/10$ on $[14,16]$; with $|x_p|<12q$,
\eqref{eq:escape} and the mean-value theorem
give
$|\kappa-3\pi^2/2|<500pq$, exactly
\eqref{eq:tip-rate}.  The strict orbit enclosure
also gives the two monotone inverse sheets and
absence of an earlier central critical return.
\end{proof}

\subsection{Proper degree-two returns (not globally degree two)}
\begin{proof}[Proof of Theorem~\ref{thm:rouche}]
For $f_c^j(z)=x_j+a_jz^2+E_j(z)$,
$a_1=1$, $a_{j+1}=2x_ja_j$, and
$E_1=0$, the \emph{identity}
\[
 E_{j+1}(z)=2x_jE_j(z)+(a_jz^2+E_j(z))^2
\]
proves by induction, \emph{on the whole complex disk} $|z|\le R$,
\begin{equation}\label{eq:tail}
 T_1=0,\quad
 T_{j+1}=2|x_j|T_j+(|a_j|R^2+T_j)^2,\quad
 \sup_{|z|\le R}|E_j(z)|\le T_j.
\end{equation}
This is a complete tail, not a finite Taylor truncation.
For $p\ge13$, $|x_j|<2$ and $|a_j|\le4^{j-1}$.
If $t_j=4^jR^2\le1$, induction yields
$T_j\le t_j^2/128$: for $j<p$,
$t_j\le1/4$ and
$4/128+(1/4+1/512)^2<16/128$.
With $R=R_p=2^{-p}$, $t_p=1$ and
$T_p\le1/128$.  Thus
\eqref{eq:A-bds}, $|x_p|<12q$ and
$R_p\le2^{-13}$ imply \eqref{eq:p-margin}.
Every one of these strict comparisons reduces
to a rational inequality at the worst $p=13$.
For $3\le p\le12$ the original ten root packets
are propagated via \eqref{eq:tail}, respecting
the exceptional $R_3=1/4$, to prove
the identical $>1/25$ boundary margin.

For any $|y|\le R_p$, on $|z|=R_p$,
\[
 |G_p(z)-y-A_pz^2|
 \le|x_p|+T_p+R_p<|A_p|R_p^2.
\]
Rouch\'e's theorem
\cite[Chapter 4, \S5.2, corollary to Theorem 18]{AhlforsComplex}
counts exactly two solutions of $G_p(z)=y$
inside the disk, with none on its boundary.
Since $G_p(0)=x_p\in U_p$ and zero has
local multiplicity two, both solutions
near this critical value lie in the central
component.  By proper continuation this
component maps onto $U_p$ with degree two.
The strict boundary inequality and
Riemann--Hurwitz
\cite[\S17.14]{Forster1981} give the
central proper quadratic-like restriction,
$V_p\Subset U_p$, with a single critical point.

For the \emph{different} normalized third return
$H=h_\tau^3$, the model $g_\tau(w)=w^2+\tau$
has the \emph{exact} degree-eight polynomial
\[
 g_\tau^3(w)=(v^2+\tau)+(4\tau v)w^2
 +(2v+4\tau^2)w^4+4\tau w^6+w^8,
 \qquad v=\tau^2+\tau.
\]
For $\tau\in J$ one has the verified bounds
$7/5<v<10/7$, $|v^2+\tau|<13/50$,
$4|\tau|v>99/10$, $2v+4\tau^2<16$,
$4|\tau|<8$.  On the \emph{whole}
complex disk $|w|\le r=3/10$,
the first two intermediate images remain in
$|w|<21/10$, so the three-step perturbation
is bounded by $1848q$.  The Rouch\'e gap is
therefore strictly larger than
\[
 \frac{99}{10}r^2-\left(
   \frac{13}{50}+16r^4+8r^6+r^8+
   1848\cdot4^{-13}+r\right)>\frac18.
\]
The third-return central local degree is two;
both intermediate model images stay away from
additional critical points.  Rouch\'e and
Riemann--Hurwitz prove the proper central
$H$ restriction.  For $3\le p\le12$ the
separate full-disk $3p$-return check applies
\eqref{eq:tail} at $R=r/|A_p|$ with outward
$2^{-192}$ tail rounding and certifies the
same $>1/8$ strict gap and intermediate
component containment.
The three phase compact real intervals are
disjoint by \eqref{eq:order1}; shrinking
their analytic neighborhoods gives a
primitive three-return restriction
inside the original $G_p$ domain.  This
shrinking proves \emph{existence} but not
a uniform modulus of the chosen neighborhoods.
\end{proof}
